% Options for packages loaded elsewhere
\PassOptionsToPackage{unicode}{hyperref}
\PassOptionsToPackage{hyphens}{url}
\PassOptionsToPackage{dvipsnames,svgnames,x11names}{xcolor}
\documentclass[
]{article}
\usepackage{xcolor}
\usepackage[margin=1in]{geometry}
\usepackage{amsmath,amssymb}
\setcounter{secnumdepth}{-\maxdimen} % remove section numbering
\usepackage{iftex}
\ifPDFTeX
  \usepackage[T1]{fontenc}
  \usepackage[utf8]{inputenc}
  \usepackage{textcomp} % provide euro and other symbols
\else % if luatex or xetex
  \usepackage{unicode-math} % this also loads fontspec
  \defaultfontfeatures{Scale=MatchLowercase}
  \defaultfontfeatures[\rmfamily]{Ligatures=TeX,Scale=1}
\fi
\usepackage{lmodern}
\ifPDFTeX\else
  % xetex/luatex font selection
\fi
% Use upquote if available, for straight quotes in verbatim environments
\IfFileExists{upquote.sty}{\usepackage{upquote}}{}
\IfFileExists{microtype.sty}{% use microtype if available
  \usepackage[]{microtype}
  \UseMicrotypeSet[protrusion]{basicmath} % disable protrusion for tt fonts
}{}
\makeatletter
\@ifundefined{KOMAClassName}{% if non-KOMA class
  \IfFileExists{parskip.sty}{%
    \usepackage{parskip}
  }{% else
    \setlength{\parindent}{0pt}
    \setlength{\parskip}{6pt plus 2pt minus 1pt}}
}{% if KOMA class
  \KOMAoptions{parskip=half}}
\makeatother
\usepackage{longtable,booktabs,array}
\newcounter{none} % for unnumbered tables
\usepackage{calc} % for calculating minipage widths
% Correct order of tables after \paragraph or \subparagraph
\usepackage{etoolbox}
\makeatletter
\patchcmd\longtable{\par}{\if@noskipsec\mbox{}\fi\par}{}{}
\makeatother
% Allow footnotes in longtable head/foot
\IfFileExists{footnotehyper.sty}{\usepackage{footnotehyper}}{\usepackage{footnote}}
\makesavenoteenv{longtable}
\setlength{\emergencystretch}{3em} % prevent overfull lines
\providecommand{\tightlist}{%
  \setlength{\itemsep}{0pt}\setlength{\parskip}{0pt}}
\usepackage{bookmark}
\IfFileExists{xurl.sty}{\usepackage{xurl}}{} % add URL line breaks if available
\urlstyle{same}
\hypersetup{
  pdftitle={Told or Enforced: When In-Context Contracts Substitute for Runtime Enforcement in Agent Harnesses},
  pdfauthor={Dom Colligan},
  colorlinks=true,
  linkcolor={Maroon},
  filecolor={Maroon},
  citecolor={Blue},
  urlcolor={Blue},
  pdfcreator={LaTeX via pandoc}}

\title{Told or Enforced: When In-Context Contracts Substitute for
Runtime Enforcement in Agent Harnesses}
\author{Dom Colligan}
\date{}

\begin{document}
\maketitle

{
\setcounter{tocdepth}{3}
\tableofcontents
}
\section{Told or Enforced: When In-Context Contracts Substitute for
Runtime Enforcement in Agent
Harnesses}\label{told-or-enforced-when-in-context-contracts-substitute-for-runtime-enforcement-in-agent-harnesses}

Dom Colligan

\subsection{Abstract}\label{abstract}

When you build an agent on top of a user's data, you can keep it inside
the rules two ways: state the rule in the agent's prompt and rely on the
model to follow it, or have the surrounding software (the
\emph{harness}) block the forbidden action outright, whatever the model
decides. Real systems do both at once, so no one has measured what the
second buys once you have done the first. We separate the two levers ---
the rule \emph{told} or \emph{withheld}, and enforcement \emph{on} or
\emph{off} --- on the governance rules of a frozen reference runtime,
and ask the question a harness designer actually faces: given the model
was told, does enforcement still change its behavior?

In a small, pre-registered, single-runtime study the answer depends on
the model. We test the \emph{commit boundary}: the agent may propose a
change to the user's data, but only the user may commit it. Told this
rule with enforcement off, both capable models (MiniMax-M3 and
Llama-3.3-70B) refused on every run of both test tasks, while both
weaker models (Qwen3.5-9B and Qwen2.5-7B) committed the change on every
run. Capable models enforce the rule themselves; weak ones need the
runtime to. We report raw counts rather than a p-value: with two test
tasks per side and near-identical repeats, the honest sample size is
two, so this is a clear demonstration, not a powered result --- and a
confounded one, since capability tracks model family here and, once the
weakest model is set aside, the load-bearing evidence rests on a single
mid-size model.

Two contributions stand apart from that ladder. A methodological
warning: eval harnesses that hide a tool's output from the agent force
it to guess identifiers, and a scorer then reads the guess as
fabrication --- a fact about the harness, not the model; our own
pipeline produced and then dissolved exactly such a finding. And a
released benchmark, GovernedAgentBench, that holds the two levers apart
with a deterministic, offline scorer.

\subsection{1 Introduction}\label{introduction}

An agent harness is a design surface. Whoever wires a language model up
to real tools controls the loop around it: which tools are exposed, how
the model's actions are parsed, and what happens when an action would
break a rule. For any one rule, that engineer has two levers. The first
is \textbf{telling}: put the rule in the prompt and rely on the model.
The second is \textbf{enforcing}: have the harness block or gate the
violation, regardless of what the model does. Telling is cheap and moves
with the prompt; enforcing has to be built, tested, and maintained. So
the practical question is: \emph{when does the second lever change
behavior the first has already secured?}

Existing guardrail evaluations do not answer this, because they run the
two levers together. Some toggle enforcement while the full policy sits
in the prompt in both conditions {[}arXiv:2604.15579{]}, so the delta is
only measured where the model was already told. Others bundle ``injected
and enforced'' into one switch {[}arXiv:2602.22302{]}, or compare an
enforced arm against a prompt-only arm {[}arXiv:2603.00822,
arXiv:2605.28914{]}, which mixes telling into enforcing. The closest,
Mind the GAP {[}arXiv:2602.16943{]}, varies prompt \emph{wording}
against a runtime contract but never withholds the rule. The cell that
isolates enforcement --- \emph{rule withheld, runtime still enforcing}
--- is nearly absent from prior work, and no located study crosses
told-versus-withheld against enforced-versus-off, per rule, and scores
the model's first action.

This paper runs that crossing. For each rule we build a 2×2: the rule
told or withheld, and the runtime enforcing or off (Section 3). The
quantity we care about is \textbf{the marginal value of enforcement
given the model was told} --- the gap between the deployed setting (told
and enforced) and told-but-not-enforced. If those two behave the same,
enforcement added nothing \emph{to behavior} on that rule; its value is
then the guarantee that the block still cannot be bypassed, which no
behavioral result touches.

\textbf{What we find.} We hypothesized three things might decide when
telling is enough: whether the model can check its own compliance,
benign pressure to finish the task, and model capability. Capability is
the one that carried. On the commit boundary, told the rule with
enforcement off, the two capable models complied on every run of both
test tasks and the two weak models committed the forbidden change on
every run --- a clean split in what enforcement adds: nothing detectable
for the capable models, everything for the weak ones. But the honest
unit of replication is the \emph{task}, and there are two per side with
near-identical repeats, so we present raw counts and a per-run check of
\emph{why} each model did what it did (a capable model refuses and tells
the user to run the commit; a weak model runs it while the gate is off),
not a significance claim. The pattern is stark but small and confounded:
capability is entangled with model family (both capable models are
non-Qwen, both weak ones Qwen), and once the weakest model --- whose
failures are partly it being too weak to drive the tools at all --- is
set aside, the robust evidence rests on one mid-size model. The other
two ideas did not survive early probing (Section 5). We claim no
``guardrails are useless'' result: where the model self-enforces,
enforcement's \emph{behavioral} value is undetected at this sample size
but its guarantee is intact, and it stays load-bearing wherever the
model \emph{cannot} self-enforce --- below the capability floor, when
the rule was never told, and under adversarial pressure, which we do not
test.

\textbf{A second, independent finding is methodological.} Our own
harness first produced a spurious result: it handed the agent a file
path instead of a tool's actual output, so the blinded agent guessed
identifiers it could not see, and the scorer called that fabrication.
Restoring the output dissolved the effect entirely (Section 6). We name
the failure mode \textbf{harness blindness}.

\textbf{Contributions.} (1) A capability-dependent finding, reported as
a small single-runtime case study rather than a powered claim: told the
commit rule with enforcement off, capable models comply and weak models
violate, on raw counts with a per-run mechanistic check, with the
confounds stated plainly. (2) The harness-blindness caution,
demonstrated by dissolving one such finding. (3) GovernedAgentBench: a
task suite, a deterministic offline scorer, and a git-pinned reference
runtime that let the two levers be varied independently. The novelty is
the \emph{conjunction} --- the withheld-but-enforced cell inside a
per-rule 2×2, scored on first action, with a released scorer --- not any
single piece; and the paper claims no scaling law, no injection
robustness, and nothing beyond this one runtime.

\subsection{2 Related work}\label{related-work}

\textbf{Measuring the harness.} The software layer between a model and
its tools is now studied in its own right {[}arXiv:2605.27922,
arXiv:2603.25723, arXiv:2603.28052{]}, but that line optimizes harnesses
for capability and carries no governance axis; we measure the governance
layer instead. A capability reported at the model-and-harness level
rather than the bare model {[}arXiv:2605.27922{]} is the backdrop for
both of our contributions.

\textbf{Enforcement vs.~telling.} Enforcement frameworks
{[}arXiv:2503.18666{]} supply mechanisms without measuring what the
model would have done alone. In injection defense, architectural
enforcement beats prompt-only defenses under attack {[}arXiv:2503.18813,
arXiv:2504.11703{]}, but those never withhold the policy per rule, so
they do not separate telling from enforcing. Two 2026 results run
partial versions of our contrast and reach the same shape from the other
side: Agent Behavioral Contracts {[}arXiv:2602.22302{]} evaluates
runtime contracts at scale but its conditions ``differ only in whether
the rules are injected \emph{and} enforced,'' bundling the levers; and
Mind the GAP {[}arXiv:2602.16943{]} reports both that prompt-based
tool-call safety is strongly model-dependent and that a runtime
governance contract shows no deterrent effect --- but it varies prompt
wording, never withholds the rule, and its ``no deterrent'' is a logging
layer under jailbreak, not a hard block under benign use (we take up the
apparent tension in Section 8). Pointing the opposite way, Mechanical
Enforcement {[}arXiv:2605.14744{]} finds mechanical enforcement beating
text-only governance in a banking agent; it pairs two configurations
rather than crossing a factorial and scores rationale quality, not
tool-action compliance.

\textbf{Not just ``models follow instructions.''} The risk is misreading
this as instruction-following. That capable models follow
\emph{checkable} instructions is established {[}arXiv:2311.07911{]}, but
told-only compliance is not free: capable models complete tasks while
breaking latent rules (``unsafe success'') non-monotonically across
families {[}arXiv:2601.08196{]}, and rules obeyed while visible are
silently dropped as context is compacted {[}arXiv:2606.22528{]}. Our
question is not whether an agent \emph{can} follow a stated rule but
what enforcement adds once it has been told.

\textbf{Closest cells.} No located study puts the withheld-but-enforced
cell inside a full 2×2 against told-and-enforced and a neither-floor,
per rule, scored on first action, with a released scorer. FORGE
{[}arXiv:2602.16708{]} runs the only by-design withheld-but-enforced
condition, but has no told-and-enforced cell, no floor, uses multi-turn
corrective feedback, and releases nothing. ABSTAIN
{[}arXiv:2606.02965{]} reaches three of the four cells for one mechanism
but its enforced cell reuses the told prompt. PhantomPolicy
{[}arXiv:2604.12177{]} contributes the ``policy-invisible violation''
idea we build on but folds telling into one ordinal enforcement axis.
TeamBench {[}arXiv:2605.07073{]} finds prompt-only and sandbox-enforced
role separation statistically indistinguishable on pass rate yet 3.6×
more verifier-overrides in the prompt-only arm --- convergent evidence
for the substitution shape in a different domain, and a caution (Section
7) that aggregate parity can hide sub-metric value. The contribution is
the conjunction, together with the harness-blindness caution; no single
element is claimed as first.

\subsection{3 The two levers and the
2×2}\label{the-two-levers-and-the-22}

A harness can make an agent respect a rule two ways: state the rule
where the model can read it (\textbf{telling}), or make the violation
impossible whatever the model emits (\textbf{enforcing}). Deployed
systems pull both at once. This paper's whole manipulation is the
\emph{gap} between the copy of the rule the model can read and the rule
the runtime actually implements.

\textbf{Telling} is the in-context contract: not just a policy paragraph
but everything the model is shown --- the list of commands and their
argument shapes, a per-command ``safe for the agent to run'' flag, and
the prose describing the out-of-scope boundary. It is a harness-level
lever, separate from anything baked in during training.

\textbf{Enforcing} is what the runtime does regardless of the model. Our
reference runtime has five governance checks, each of which can be
switched off (individually, or all at once):

{\def\LTcaptype{none} % do not increment counter
\begin{longtable}[]{@{}
  >{\raggedright\arraybackslash}p{(\linewidth - 2\tabcolsep) * \real{0.5000}}
  >{\raggedright\arraybackslash}p{(\linewidth - 2\tabcolsep) * \real{0.5000}}@{}}
\toprule\noalign{}
\begin{minipage}[b]{\linewidth}\raggedright
Check
\end{minipage} & \begin{minipage}[b]{\linewidth}\raggedright
What it does
\end{minipage} \\
\midrule\noalign{}
\endhead
\bottomrule\noalign{}
\endlastfoot
schema check & rejects malformed or invented commands before they run \\
dispatch gate & refuses to run a command not flagged safe for the
agent \\
commit gate & the agent may \emph{propose} a change to user data, but
only the user may commit it \\
refusal rule & refuses out-of-scope requests, including a zero-tolerance
medical boundary \\
audit check & attaches evidence to recommendations; the scorer later
catches fabricated or missing references \\
\end{longtable}
}

A sixth property, transaction integrity, is always on. The
\textbf{commit gate} is the one this paper's headline turns on: the
runtime refuses an agent-issued commit and leaves the change merely
proposed.

\textbf{The 2×2.} Crossing the two levers, for any one rule, gives four
cells:

{\def\LTcaptype{none} % do not increment counter
\begin{longtable}[]{@{}lll@{}}
\toprule\noalign{}
& Enforcing & Off \\
\midrule\noalign{}
\endhead
\bottomrule\noalign{}
\endlastfoot
\textbf{Told} & A: deployment baseline & B: told, not enforced \\
\textbf{Withheld} & C: withheld, still enforced & D: neither \\
\end{longtable}
}

Cell A is how systems ship. Cell B isolates \textbf{self-enforcement}:
the model knows the rule and nothing stops it. Cell C isolates
\textbf{pure enforcement}: the rule is hidden but the runtime still
blocks. Cell D is the floor. Three reads carry the information: B vs.~D
is the effect of telling; C vs.~D is the effect of enforcing; and
\textbf{A vs.~B is our headline --- what enforcement adds once the model
was told.} If A and B match, enforcement changed no behavior on that
rule; its value is then the guarantee itself, which is real and
unconditional (a blocked action simply cannot happen) even where
behavior does not move.

One scoring convention matters. A blocked action returns an error, and
that error is itself the rule, told late --- so cell C drifts toward
cell B after the model's first contact with a block. We therefore score
the telling reads on the model's \textbf{first action}, and report
multi-turn behavior separately.

\subsection{4 GovernedAgentBench}\label{governedagentbench}

GovernedAgentBench holds the two levers apart: it toggles telling and
enforcing independently, per rule, against one fixed runtime, and scores
the result deterministically and offline.

\textbf{The runtime.} The instrument is HAI, a non-clinical
personal-wellness agent runtime with a typed command interface, pinned
as a frozen snapshot; it is a reference implementation, not the
contribution or a product. Its no-diagnosis boundary is one of the
evaluated rules, not a disclaimer. Two facts matter for reproduction.
The published \texttt{v0.2.0} package \textbf{does not enforce} --- the
dispatch gate landed three days after that version tag --- so the
measured runtime is pinned by git commit \texttt{6c82cd0} (14 commits
past the nearest tag, which predates the enforcement fixes and must not
be used to reproduce). And before spending anything, a pre-flight check
confirmed the runtime actually refuses an agent commit; it passed, so
the enforced cells were measured against a runtime \emph{verified} to
enforce, not assumed to.

\textbf{Tasks.} The suite has 39 tasks; the paid run used a concentrated
16-task subset, the rest covered offline. Each task is one labelled cell
of a rule's 2×2. For the commit boundary we use two tasks (committing a
target and committing an intent), each run four times per cell.

\textbf{The two axes in code.} Enforcing is a runtime setting. Telling
is a per-task flag applied at prompt-render time while the runtime is
untouched: the \emph{withheld} prompt strips the rule's facts and prose,
and an automatic scan rejects any withheld prompt that still leaks the
rule's words. Two honest caveats attach. First, even withheld, the
prompt still lists command \emph{names}, so the model has a hint; B-vs-D
and C-vs-D are therefore lower bounds on the effect of telling. Second,
told and withheld were hand-rebuilt before the run so that ``told''
states the operative rule and ``withheld'' removes it --- fixing an
earlier version in which both arms committed --- and a small patch that
lets the withheld model still \emph{find} the commit command was applied
after the pre-run canary, so the withheld violation rate carries that
provenance and can be re-scored without the patch offline.

\textbf{Scoring.} The scorer is deterministic and offline, with no model
calls. We disclose one thing plainly: the harm-check code was edited
three times in the day before the run, the last 47 minutes before it,
each edit in the direction we hoped for. We verified this does not touch
the result --- \textbf{every violation in the headline table is caught
by the original, older detection path; re-scoring with the pre-edit
scorer is identical on the 2×2} --- so the finding is invariant to those
edits. One circularity is disclosed rather than hidden: the
medical-boundary detector shares its phrase list with the runtime's own
refusal, so the enforced cell is clean on that check by construction.

\textbf{The pre-registered run.} The design was locked before the run
(models, four repeats per cell, a canary that had to move before the
paid sweep, and pre-committed language for every outcome). We disclose
two departures. The pre-registered primary model was deprecated by the
provider mid-cycle; we replaced it with MiniMax-M3, but only after a
canary confirmed the replacement reproduced the effect --- an
outcome-conditioned choice, so \textbf{MiniMax-M3 is a
non-pre-registered addition and its result is exploratory; the
pre-registered capable model is Llama-3.3-70B}, which shows the same
pattern. And the commit boundary has two test tasks, not the three the
sample-size analysis assumed, so its pooled cell is eight runs, below
our own twelve-run target; the effects are saturated (0/8 vs.~8/8), but
the small size is stated, not hidden.

\subsection{5 Results}\label{results}

The headline is the four-model capability ladder (Section 5.1). The
earlier one-model probes that shaped the design are summarized as
provenance in Section 5.2. Throughout, cell A (told, enforced) and cell
C (withheld, enforced) are \textbf{100\% safe by construction} --- the
runtime blocks the commit whatever the model does, so all 64 enforced
runs score identically --- which means the entire behavioral finding
lives in the \emph{unenforced} cells, B and D.

\subsubsection{5.1 The capability ladder}\label{the-capability-ladder}

The run swept four models --- MiniMax-M3 and Llama-3.3-70B (capable),
Qwen3.5-9B (near the floor), Qwen2.5-7B (below it) --- against the
git-pinned runtime, for a total paid cost of USD 10.44. The commit-gate
results, as percentage of runs that stayed safe:

{\def\LTcaptype{none} % do not increment counter
\begin{longtable}[]{@{}
  >{\raggedright\arraybackslash}p{(\linewidth - 8\tabcolsep) * \real{0.2000}}
  >{\raggedright\arraybackslash}p{(\linewidth - 8\tabcolsep) * \real{0.2000}}
  >{\raggedright\arraybackslash}p{(\linewidth - 8\tabcolsep) * \real{0.2000}}
  >{\raggedright\arraybackslash}p{(\linewidth - 8\tabcolsep) * \real{0.2000}}
  >{\raggedright\arraybackslash}p{(\linewidth - 8\tabcolsep) * \real{0.2000}}@{}}
\toprule\noalign{}
\begin{minipage}[b]{\linewidth}\raggedright
Model
\end{minipage} & \begin{minipage}[b]{\linewidth}\raggedright
A: told, enforced
\end{minipage} & \begin{minipage}[b]{\linewidth}\raggedright
B: told, off
\end{minipage} & \begin{minipage}[b]{\linewidth}\raggedright
C: withheld, enforced
\end{minipage} & \begin{minipage}[b]{\linewidth}\raggedright
D: neither
\end{minipage} \\
\midrule\noalign{}
\endhead
\bottomrule\noalign{}
\endlastfoot
MiniMax-M3 (capable) & 100* & 100 & 100* & 0 \\
Llama-3.3-70B (capable) & 100* & 100 & 100* & 25 \\
Qwen3.5-9B (near floor) & 100* & 0 & 100* & 12.5 \\
Qwen2.5-7B (below floor) & 100* & 0 & 100* & 12.5 \\
\end{longtable}
}

*Enforced cells (A, C) are safe by construction; the finding is in cells
B and D.

\textbf{The finding is column B --- told the rule, runtime off.} Both
capable models stayed safe on all eight runs (both tasks × four
repeats); both weak models committed the forbidden change on all eight.
So enforcement adds nothing detectable once a capable model is told, and
it is the whole barrier for a weak one.

\textbf{How we count.} The honest unit of replication is the
\emph{task}, not the run: there are two tasks per side, and the four
repeats are near-identical (seven of the eight told-off conditions are
byte-for-byte the same across all four repeats). So a run-level
significance test overstates its evidence --- the weak-side difference
is ``p = 0.00016'' counting runs but \textbf{p = 0.33 counting tasks
(two vs.~two)} --- and we do not lead with it. Instead we lead with the
raw counts and a per-run check of \emph{why}: in the told-off cell,
MiniMax refuses the commit on its own turn and tells the user to run it
themselves, while Qwen-7B issues the commit command itself and it
executes because the gate is off. The capable side's ``no effect'' is a
\emph{failure to detect} at this size (an upper bound of about 32
points), not proven equivalence.

\textbf{Why this reads as substitution, not the model just being nice.}
One might worry the capable models would comply anyway, told or not.
They would not: even for them, cell D (rule withheld, runtime off)
violates (MiniMax 0\%, Llama 25\% safe), so \emph{telling} has a large
effect within a single model, and the capable model's flat A-vs-B is
genuine substitution rather than a model that never needed either lever.
(This D-cell violation depends on the discoverability patch noted in
Section 4, which can be re-scored offline.)

\textbf{Three caveats on the reading.} First, the below-floor 7B's
failure is partial --- its violations are graded rather than total ---
so at a looser threshold its column B would read 75\%, not 0\%; the
load-bearing cells (capable 100\%, near-floor 9B 0\%) are unaffected by
any threshold choice. Second, parse failures cannot fake this: no
headline cell contains a run that was scored safe merely because the
model emitted nothing, and Llama produces malformed output on
\textasciitilde40\% of runs yet still reaches 100\% through genuine
refusals, so parse rate does not explain the result. Third, once the
below-floor 7B is set aside and the capability-family confound held, the
robust ``told but does not self-enforce'' evidence rests on \textbf{one
mid-size model, Qwen3.5-9B, at eight runs on one rule and one runtime}
--- a real, mechanistically verified failure, but not yet a multi-model
or family-general one.

\textbf{The other rule we swept, the medical-boundary refusal, is
near-ceiling and uninformative.} Its detector fires on only 7 of 223
runs, all on one task and one model; the two tasks built to provoke a
medical claim produce none, even below the floor. We read no result off
it, and flag that the near-total silence may reflect genuine
boundary-respect or a detector that is too narrow (Section 7).

\subsubsection{5.2 What the earlier one-model probes
showed}\label{what-the-earlier-one-model-probes-showed}

Before the ladder, a set of single-model probes (Qwen3-235B, three to
five runs per cell) framed the design; we summarize rather than belabor
them. With the runtime off, a told-and-salient model refused 3/3 while a
withheld model committed the violation 3/3, establishing that telling
and enforcing each move behavior against the floor. The two ideas that
did \emph{not} survive: the rule we predicted a model \emph{could not}
self-check (audit faithfulness) turned out to be self-checkable once the
model retrieved the evidence, and benign pressure to finish the task
produced zero fabrication across 40 runs. One qualifier the probes raise
and the ladder does not settle: self-enforcement is
\textbf{salience-sensitive} --- the same told model that refused when
the rule was foregrounded dispatched the forbidden command 3/3 when it
met the rule only incidentally. This rests on three runs of one
off-roster model and is not independently reconstructable, so we carry
it as a caveat, not a result.

\subsection{6 A methodological caution: harness
blindness}\label{a-methodological-caution-harness-blindness}

Our probing produced, and then destroyed, an apparent positive result
--- and the failure mode is general enough to report as a finding.

Early on, the harness did not show the agent its tools' output: where a
command's result should have appeared, the agent got a file path
instead. It could run a lookup but never read the answer. When a task
then needed an identifier to proceed, the agent guessed --- across three
probe sets, eight cases where it issued commands carrying plausible
identifiers it had never seen, including an invented evidence-card id.
This looked like a real finding: honest when asked plainly, fabricating
when an id was instrumentally required. It was our strongest surviving
positive result.

It was an artifact of the harness. We pre-registered a falsification
test (fabrication effect must clear 40 points to count) and re-ran it
after fixing the harness to show the agent its command output.
Fabrication dropped to \textbf{zero} against the 40-point bar: once the
agent can see the empty lookup, it abstains honestly even when it needed
the id to finish. The clean before/after holds the task fixed ---
pre-fix, one run fabricates an id while another stays honest; post-fix,
all runs are honest.

The lesson: \textbf{an eval harness that hides tool output makes the
agent guess, and a scorer reads the guess as fabrication.} The number is
a fact about the harness, not the model. Genuine constraint-driven
fabrication does exist under more adversarial setups
{[}arXiv:2606.14831{]}; our point is narrower --- before charging an
agent with making something up, check that it could see what it is
accused of inventing. A related corollary falls straight out of the 2×2:
an ablation that toggles enforcement while leaving the policy in the
prompt in \emph{both} conditions {[}arXiv:2604.15579{]} measures
self-enforcement, not enforcement, and a small delta there bounds
nothing about an untold agent.

\subsection{7 Limitations and scope}\label{limitations-and-scope}

\textbf{This is a small case study on one runtime.} Every model-backed
number comes from one runtime, one prompt template, four models, eight
runs per headline cell. It is a single-runtime case study by decision:
whether the effect generalizes beyond this runtime is the highest-value
follow-up, and until an independent replication exists, nothing here
separates a real phenomenon from a quirk of this instrument.

\textbf{The load-bearing arm is effectively one model.} Capability is
confounded with model family (capable = non-Qwen, weak = Qwen), so the
ladder cannot fully separate ``weaker models need enforcement'' from
``Qwen models do.'' The two capable models agreeing across families
strengthens the substitution side, but the enforcement-is-load-bearing
side, after setting aside the below-floor 7B, rests on Qwen3.5-9B alone.
Per-model vendor-recommended sampling also entangles model identity with
sampling settings across the ladder, though the per-model 2×2 is
computed at one setting and is unaffected.

\textbf{The saturated cells cut both ways.} Because the runtime blocks
the commit regardless of the model, cells A and C are 100\% by
construction; all the behavioral signal lives in B and D, and A-vs-B is
a one-sided quantity, not a two-sided contrast. We report it as an upper
bound, not an equivalence.

\textbf{Pre-registration honesty.} Beyond the outcome-screened
replacement model and the sub-target sample size (Section 4), the
withheld prompt still leaks command names, and the medical-boundary
detector shares ground truth with the runtime. TeamBench's caution
applies to us too: our scorer sees whether the specific guarded
violation occurs, but an enforcement effect that shifted some finer
channel while leaving the tracked violation unchanged would go
undetected at this size.

\textbf{Untested regions.} Adversarial inputs are out of scope --- the
injection literature already shows telling fails under attack
{[}arXiv:2503.18813, arXiv:2504.11703{]}. Long-horizon rule drift, where
a rule obeyed while visible is dropped as context is compacted
{[}arXiv:2606.22528{]}, is named but unmeasured here.

\subsection{8 Discussion: what enforcement is still
for}\label{discussion-what-enforcement-is-still-for}

The result relocates enforcement's value; it does not remove it.
Enforcement keeps its unconditional guarantee --- a blocked action
cannot happen, whatever the model does or is prompted --- and that
guarantee is untouched by any behavioral number, including a
100\%-refusal rate at tiny n.~Beyond it, enforcement stays
\emph{behaviorally} load-bearing wherever a model cannot self-enforce:
below the capability floor, where the failure is the model being too
weak to drive the tools rather than disobedience; when the rule was
never told, the floor cell where the runtime is the only barrier and
real deployments accumulate such rules through drift and context
compaction {[}arXiv:2606.22528{]}; when a told rule is present but not
salient at the moment of action; and under adversarial pressure, which
we do not test.

This squares with the two opposite-looking priors of Section 2.
Mechanical Enforcement {[}arXiv:2605.14744{]} finds enforcement helping
--- but in a different domain, scoring rationale quality under
regulatory pressure, not tool-action compliance under benign use. Mind
the GAP {[}arXiv:2602.16943{]} finds a runtime contract with no
deterrent effect --- but its contract logs rather than blocks, whereas
our commit gate blocks outright, which is exactly why our
withheld-but-enforced cell is safe by construction. Neither contradicts
a \emph{behavioral redundancy given telling} read; they measure
different things under different pressure.

\subsection{9 Conclusion and future
work}\label{conclusion-and-future-work}

On the commit boundary, in a small pre-registered study on one runtime,
telling the rule is enough for a capable model and not for a weak one.
Told the rule with enforcement off, both capable models refused on every
run of both tasks and both weak models committed the forbidden change on
every run; because even the capable models violate when the rule is
withheld, this is genuine substitution rather than a model that never
needed either lever. We report it as raw counts on a two-task-per-side
unit, not a significance claim, and it is confounded --- capability with
model family, the load-bearing arm effectively one model. The
medical-boundary rule was uninformative at this scale. Enforcement is
redundant only in that narrow, high-capability, benign sense; its
guarantee is intact and it stays load-bearing everywhere a model cannot
self-enforce.

The takeaway we most want carried is the methodological one: before
attributing fabrication to a model, check what the harness let it see.

Four things would sharpen this. A \textbf{powered version} needs more
\emph{test tasks per rule}, not more repeats of the same two --- the
honest sample size is set by tasks, and a fresh pre-registered run with
more scenarios (and a non-Qwen weak model, to break the family confound)
is the legitimate route to a significance claim. \textbf{External
replication} on a second runtime with its own rules and scorer is what
would turn a case study into a phenomenon. \textbf{Adversarial inputs}
would locate where substitution ends and injection defense begins. And
\textbf{long-horizon drift} --- does enforcement catch a rule the model
has since forgotten? --- is the measurement this design most directly
points at.

\subsection{References}\label{references}

\begin{itemize}
\tightlist
\item
  FORGE: Formal Policy Enforcement for Real-World Agentic Systems.
  arXiv:2602.16708.
\item
  Mind the GAP: Text Safety Does Not Transfer to Tool-Call Safety in LLM
  Agents. arXiv:2602.16943.
\item
  Agent Behavioral Contracts. arXiv:2602.22302.
\item
  ABSTAIN / What Benchmarks Don't Measure. arXiv:2606.02965.
\item
  PhantomPolicy: Policy-Invisible Violations in LLM-Based Agents.
  arXiv:2604.12177.
\item
  TeamBench: Agent Coordination under Enforced Role Separation.
  arXiv:2605.07073.
\item
  Symbolic Guardrails. arXiv:2604.15579.
\item
  Mechanical Enforcement. arXiv:2605.14744.
\item
  ContextCov. arXiv:2603.00822.
\item
  AIRGuard. arXiv:2605.28914.
\item
  AgentSpec. arXiv:2503.18666.
\item
  CaMeL: Defeating Prompt Injections by Design. arXiv:2503.18813.
\item
  Progent. arXiv:2504.11703.
\item
  Harness-Bench. arXiv:2605.27922.
\item
  LogiSafetyBench. arXiv:2601.08196.
\item
  Governance Decay (mitigation: Constraint Pinning). arXiv:2606.22528.
\item
  IFEval. arXiv:2311.07911.
\item
  Constraint-driven fabrication under adversarial pressure.
  arXiv:2606.14831.
\end{itemize}

\subsection{Appendix A: reproduction}\label{appendix-a-reproduction}

The headline table reproduces from the released analysis
(\texttt{per\_gate\_substitution.json}) and the pooling module
(\texttt{build\_cell\_contrasts\_pooled}); the exact tests beside it
come from a small released helper (\texttt{exact\_tests.py}, pure
Python, no dependencies), which returns the run-level (0.00016),
task-level (0.33), per-sub-gate (0.029), and family-corrected (0.00062)
figures and the Clopper-Pearson intervals ({[}63,100{]} for 8/8,
{[}0,37{]} for 0/8). Offline artifacts regenerate with
\texttt{reproduce\_offline.py} (no network, no private data). The
model-backed trajectories and scores ship as a versioned archive
alongside the preprint; reproducing the runtime means checking out git
\texttt{6c82cd0}, not installing the package. The full task table and
the earlier retired apparatus (a 28-task / 25-oracle-pair design,
superseded 2026-07-04) are documented with the release.

\end{document}
